% Estilo de los diagramas de clases de la API (clases_*.tex). Va DESPUÉS de
% ../../figuras/estilo_tikz.tex: usa su paleta y sus puntas de flecha.
%
% Notación: la de UML, recortada a lo que hace falta acá.
%   caja de tres partes        nombre · atributos · métodos
%   triángulo hueco, llena     hereda (en C: el struct base como primer miembro)
%   triángulo hueco, punteada  cumple la interfaz sin heredar (duck typing)
%   rombo lleno                compone: el dueño crea la parte y la cierra
%   flecha punteada            usa, o recibe por parámetro
% El color del encabezado dice qué clase de clase es; va en la leyenda de cada
% figura.
%
% Los nombres llevan '_' sin escapar: cada figura pone \catcode`\_=12 al
% abrir el documento.

\tikzset{
  clase/.style      = {rectangle split, rectangle split parts=3,
                       rectangle split part align={center,left,left},
                       rectangle split part fill={azul!30,white,white},
                       draw=black, line width=0.6pt, inner sep=3pt,
                       text width=#1, font=\fontsize{6.2}{7.5}\selectfont\ttfamily},
  clase/.default    = 4.4cm,
  falsa/.style      = {clase=#1, rectangle split part fill={grisr,white,white}},
  falsa/.default    = 4.4cm,
  shim/.style       = {clase=#1, rectangle split part fill={amarillo,white,white}},
  shim/.default     = 4.4cm,
  interfaz/.style   = {clase=#1, rectangle split part fill={verde!45,white,white},
                       dashed},
  interfaz/.default = 4.4cm,
  hilo/.style       = {clase=#1, rectangle split part fill={naranja!70,white,white}},
  hilo/.default     = 4.4cm,
  % un tipo de C (libmca): struct con sus campos y las funciones que lo reciben
  ctipo/.style      = {clase=#1, rectangle split part fill={rojo!30,white,white}},
  ctipo/.default    = 4.4cm,
  % caja de una sola parte: módulos de funciones, referencias a otra figura
  modulo/.style     = {draw=black, line width=0.6pt, inner sep=3pt, align=left,
                       text width=#1, fill=white,
                       font=\fontsize{6.2}{7.5}\selectfont\ttfamily},
  modulo/.default   = 4.4cm,
  hereda/.style     = {-{Triangle[open,length=2.6mm,width=2.8mm]}, line width=0.6pt},
  cumple/.style     = {hereda, dashed},
  compone/.style    = {{Diamond[length=3.2mm,width=2.2mm]}->, line width=0.6pt},
  usa/.style        = {->, densely dashed, line width=0.55pt, black!70},
  rama/.style       = {line width=0.6pt},
  rol/.style        = {font=\fontsize{5.8}{6.8}\selectfont\ttfamily, inner sep=1.2pt,
                       fill=white, text=black!80},
}

% encabezado: nombre de la clase y, debajo, de dónde sale
\newcommand{\clase}[2]{{\bfseries\fontsize{7.8}{9.2}\selectfont #1}\\[0.8pt]%
  {\fontsize{5.6}{6.6}\selectfont #2}}
% estereotipo: «classmethod», «prop», «static», «opcional»
\newcommand{\es}[1]{{\normalfont\itshape\fontsize{5.4}{6.4}\selectfont «#1»}}
% rótulo de grupo dentro de la caja de métodos
\newcommand{\gr}[1]{{\normalfont\itshape\color{black!55}#1}\\}

% La misma leyenda que \leyenda de estilo_tikz.tex, con el ancho de la nota
% como argumento: la común lo fija en 19 cm y eso ensancha la figura entera.
% #1 = posición, #2 = ancho de la nota, #3 = color/texto, #4 = nota.
\newcommand{\leyendac}[4]{%
  \begin{scope}[shift={(#1)}, every node/.style={font=\scriptsize}]
    \coordinate (t0) at (-0.5,0);
    \foreach \c/\t [count=\i] in {#3} {
      \pgfmathtruncatemacro\j{\i-1}
      \node[draw, fill=\c, minimum width=4mm, minimum height=3mm, inner sep=0pt,
            right=5mm of t\j.east, anchor=west] (l\i) {};
      \node[right=1.5mm of l\i] (t\i) {\t};
    }
    \node[anchor=north west, black!70, text width=#2, align=left]
      at (-0.05,-0.3) {#4};
  \end{scope}}
